\section{Introduction}

Category theory has often been regarded as the mathematics of mathematics~\cite{maclane1998categories}.
The capacity to lift particular concepts into a more general and abstract space
has demonstrated, time and again, that
all of mathematics seems to be governed by universal constructions.
The investigation of such universal concepts allows us to see what a particular view
could not, and to realize the intrinsic connections and shared nature of what was once thought
to be separate.

One such field, that appears throughout mathematics
yet stands apart, is optimization.
Optimal thinking is present in every mathematical construct in which one seeks the
``ideal'' object in some criterion. This notion of ``ideal'' often refers to the
simplest object that meets certain criteria. Prior work has already characterised the internal
optimization structure within categorical constructs \cite{Khadh2026bridging},
but such results are still isolated; a unified treatment is still to be formalized.

This work is motivated by the observation that many important
mathematical structures used in optimization can be understood within
a single compositional framework \cite{stein2024compositional, hanks2024compositionalopt}.
Instead of treating relations, convex functions, polyhedra,
and linear programs as isolated objects, we show that they can all be
interpreted as morphisms in categories whose composition is given by elimination or minimization.
In this way, optimization itself becomes an algebraic process and a universal setting across mathematics.

This expressivity becomes especially apparent when analyzed through
the language of string diagrams \cite{joyal1991geometry, selinger2010survey, piedeleu2023introductionstringdiagramscomputer}.
Originating as a rigorous graphical language for
monoidal categories, string diagrams provide an intuitive yet mathematically
precise syntax in which the sequential and parallel composition of morphisms
correspond directly to vertical and horizontal juxtaposition of diagrams
\cite{joyal1991geometry, maclane1998categories}. Building on this foundation,
a rich family of \emph{diagrammatic} or \emph{graphical calculi} has emerged,
offering complete axiomatisations of various algebraic structures purely through
equational diagram rewriting \cite{selinger2010survey, bonchi2017interacting}.
These calculi have proved remarkably versatile, finding applications across a
wide range of fields: in \emph{quantum computing} \cite{coecke2011interacting};
\emph{electronics} \cite{baez2018compositional}; \emph{linear algebra} \cite{bonchi2017interacting,
	bonchi2019affine, defreitas2025generalizinginvertiblematrixtheorem};
and in \emph{optimisation and probability} \cite{stein2024gqa,stein2025compositional,
	bonchi2021polyhedral}. Throughout this work
we explore how, step by step, one can enrich the already known theories for linear
\cite{bonchi2017interacting} algebra by adding one new generator at a time, along with a particular
set of axioms, and derive even richer structures for convex problems \cite{stein2024compositional}.

A first instance of this mechanism is the Graphical Quadratic Algebra (\textbf{GQA}) \cite{stein2024gqa},
where the addition of a single ``measurement'' generator expands the expressivity of the previous theories
by acting as a sort of measurement device along the diagram, lifting diagrams from the class of binary relations
to the broader class of weighted relations. From the optimization perspective, this is what plays the role of
the objective function in an optimization problem. We do not develop the quadratic branch here, but it serves
as the conceptual template for the cost generator at the heart of our contribution.

The polyhedral branch is the one our work builds on directly. Extending the Symmetric Monoidal Theory of
Graphical Affine Algebra (\textbf{GAA}) \cite{bonchi2019affine} with an inequality generator $\tikz[baseline=-0.5ex]\pic{ge};$ yields
Graphical Polyhedral Algebra (\textbf{GPA}) \cite{bonchi2021polyhedral}. Geometrically,
this allows the algebra to go beyond affine spaces and start representing polyhedral figures as linear inequality constraints in optimization,
\[
	\{| x \in P\} =
	\begin{cases}
		0 \quad \text{if} \;Ax +b\ge 0 \\
		+\infty \quad \text{otherwise.}
	\end{cases}
\]
Inspired by these observations, we make our central contribution: the design of a graphical
algebra for parametric linear programming through the introduction of 
a the new generator $\tikz \pic{linear};$ and three axioms. In this
setting, a linear program is no longer viewed only as
a numerical problem to be solved, but as a morphism that
can be composed with other optimization blocks.
This gives a symbolic language in which the elimination of internal variables
corresponds directly to categorical composition, allowing optimization problems
to be manipulated diagrammatically before computation.

In particular, every piecewise affine convex function can be represented as the value function of a linear program \cite{rockafellar1970convex}, which means the Graphical Linear Programming (\textbf{GLP}) theory provides a syntax for a broad class of nonsmooth convex models that appear in optimization, control, economics, and machine learning~\cite{boyd2004convex} such as Feed Forward Networks architectures, eletrical circuits and others. In Section~\ref{sec:glp} we show the diagrammatic equivalence of these two classes of objects and how diagrammatically they are not only denotationally equivalent but in fact instances of the same normal form. Thus, our diagrammatic syntax blurs the distinction between expressing LP programs as optimization problems and as piecewise convex functions, deriving a canonical normal form.


More broadly, this work suggests that graphical methods can play for optimization the same role that graphical linear algebra plays for matrices: providing normal forms, symbolic rewriting rules, and completeness results within a rigorous categorical setting. This opens the possibility of using diagrammatic reasoning not only for proofs, but also for designing compositional optimization systems and formal symbolic calculi for convex programming.

\paragraph{Overview}
Section~\ref{sec:bg} introduces the algebraic and categorical playground on which the
rest of the work is built: the initial notions of quantales and relations, the Extended-Lawvere
quantale of costs, the language of symmetric monoidal theories, and the prior graphical theories
$\mathbf{GLA}$, $\mathbf{GAA}$ and $\mathbf{GPA}$ on which our contribution rests.
Section~\ref{sec:sem} introduces, for the first time, the categorical semantics of parametric linear
programs, establishing the equivalence between piecewise affine convex functions, polyhedral
functions and parametric linear programs, and packaging them as the categories $\mathbf{LPRel}$ and
$\mathbf{PolyFun}$. Section~\ref{sec:glp} defines the symmetric monoidal theory $\mathbf{GLP}$,
its normal form, and the epigraph form. Section~\ref{sec:complete} proves soundness, definability
and completeness, establishing that $\mathbf{GLP}$ axiomatises the semantics.
This work is a step toward a complete diagrammatic calculus for convex optimization.

